% GENERATED FILE. Edit canonical lessons, then run npm run generate:books.
% canonical-source-hash: fnv1a64:87c373a7a5c1ef50
% canonical-lessons: GE-C209-buerger, GE-C209-wahl, GE-C209-partei, GE-C209-politik, GE-C209-wirtschaft

\chapter{A citizen, An election, A party, Politics, The economy}
\label{ch:ge-a-citizen-an-election-a-party-politics-the-economy}
\hlchaptermodality{209}

\begin{chapteropening}
\textbf{By the end of this chapter:} \emph{I can say a citizen, an election, a party, politics and the economy.}

Complete the last lesson of chapter 209: I can say a citizen, an election, a party, politics and the economy.
\end{chapteropening}

\section[der Bürger]{der Bürger — a citizen}

\label{lesson:GE-C209-buerger}



\begin{quote}
Before the new one: say the German for a state, then the German for the population.
\end{quote}

\subsection*{You'll want to know: der Bürger}
\textbf{der Bürger} — \textquotedblleft{}a citizen\textquotedblright{}.

\begin{grammarlens}[title={What you've built}]
One more word for this chapter.
\end{grammarlens}

\subsection*{Your turn}
\begin{itemize}
\raggedright
  \item \emph{Say it:} \emph{der Bürger}
  \item \emph{Say it:} \emph{der Bürger}, once more
  \item \emph{Say it:} say \emph{der Bürger}
  \item \emph{Recall:} say the German for a state, then the German for the population, then say \emph{der Bürger} again
\end{itemize}

\begin{tcolorbox}[breakable,colback=teal!4,colframe=teal!35!black,title={Before you move on}]
What does der Bürger mean? (\textquotedblleft{}A citizen\textquotedblright{}.) Say it once more.
\end{tcolorbox}
\section[die Wahl]{die Wahl — an election, a choice}

\label{lesson:GE-C209-wahl}



\begin{quote}
Before the new one: say the German for the population, then the German for a citizen.
\end{quote}

\subsection*{You'll want to know: die Wahl}
\textbf{die Wahl} — \textquotedblleft{}an election, a choice\textquotedblright{}.

\begin{grammarlens}[title={What you've built}]
One more word for this chapter.
\end{grammarlens}

\subsection*{Your turn}
\begin{itemize}
\raggedright
  \item \emph{Say it:} \emph{die Wahl}
  \item \emph{Say it:} \emph{die Wahl}, once more
  \item \emph{Say it:} say \emph{die Wahl}
  \item \emph{Recall:} say the German for the population, then the German for a citizen, then say \emph{die Wahl} again
\end{itemize}

\begin{tcolorbox}[breakable,colback=teal!4,colframe=teal!35!black,title={Before you move on}]
What does die Wahl mean? (\textquotedblleft{}An election, a choice\textquotedblright{}.) Say it once more.
\end{tcolorbox}
\section[die Partei]{die Partei — a party (political)}

\label{lesson:GE-C209-partei}



\begin{quote}
Before the new one: say the German for a citizen, then the German for an election.
\end{quote}

\subsection*{You'll want to know: die Partei}
\textbf{die Partei} — \textquotedblleft{}a party (political)\textquotedblright{}.

\begin{grammarlens}[title={What you've built}]
One more word for this chapter.
\end{grammarlens}

\subsection*{Your turn}
\begin{itemize}
\raggedright
  \item \emph{Say it:} \emph{die Partei}
  \item \emph{Say it:} \emph{die Partei}, once more
  \item \emph{Say it:} say \emph{die Partei}
  \item \emph{Recall:} say the German for a citizen, then the German for an election, then say \emph{die Partei} again
\end{itemize}

\begin{tcolorbox}[breakable,colback=teal!4,colframe=teal!35!black,title={Before you move on}]
What does die Partei mean? (\textquotedblleft{}A party (political)\textquotedblright{}.) Say it once more.
\end{tcolorbox}
\section[die Politik]{die Politik — politics}

\label{lesson:GE-C209-politik}



\begin{quote}
Before the new one: say the German for an election, then the German for a party.
\end{quote}

\subsection*{You'll want to know: die Politik}
\textbf{die Politik} — \textquotedblleft{}politics\textquotedblright{}.

\begin{grammarlens}[title={What you've built}]
One more word for this chapter.
\end{grammarlens}

\subsection*{Your turn}
\begin{itemize}
\raggedright
  \item \emph{Say it:} \emph{die Politik}
  \item \emph{Say it:} \emph{die Politik}, once more
  \item \emph{Say it:} say \emph{die Politik}
  \item \emph{Recall:} say the German for an election, then the German for a party, then say \emph{die Politik} again
\end{itemize}

\begin{tcolorbox}[breakable,colback=teal!4,colframe=teal!35!black,title={Before you move on}]
What does die Politik mean? (\textquotedblleft{}Politics\textquotedblright{}.) Say it once more.
\end{tcolorbox}
\section[die Wirtschaft]{die Wirtschaft — the economy}

\label{lesson:GE-C209-wirtschaft}



\begin{quote}
Before the new one: say the German for a party, then the German for politics.
\end{quote}

\subsection*{You'll want to know: die Wirtschaft}
\textbf{die Wirtschaft} — \textquotedblleft{}the economy\textquotedblright{}.

\begin{grammarlens}[title={What you've built}]
One more word for this chapter.
\end{grammarlens}

\subsection*{Your turn}
\begin{itemize}
\raggedright
  \item \emph{Say it:} \emph{die Wirtschaft}
  \item \emph{Say it:} \emph{die Wirtschaft}, once more
  \item \emph{Say it:} say \emph{die Wirtschaft}
  \item \emph{Recall:} say the German for a party, then the German for politics, then say \emph{die Wirtschaft} again
\end{itemize}

\begin{tcolorbox}[breakable,colback=teal!4,colframe=teal!35!black,title={Before you move on}]
What does die Wirtschaft mean? (\textquotedblleft{}The economy\textquotedblright{}.) Say it once more.
\end{tcolorbox}
